\documentclass[10pt,a4paper]{amsart}
\usepackage[margin=19mm]{geometry}
\usepackage{amsmath,amssymb,mathtools}
\usepackage[scr=rsfs,cal=euler]{mathalfa}
\usepackage{xfrac}
\usepackage[unicode,hidelinks,bookmarks=false]{hyperref}
\newcommand{\ie}{i.e.\ }
\newcommand{\cf}{cf.\ }
\newcommand{\resp}{respectively\ }
\newcommand{\Subsection}{\subsection}
\providecommand{\nobreakdash}{\nobreak\hbox{-}}
\setlength{\emergencystretch}{2em}
\multlinegap0pt
\usepackage{amssymb}
\newtheorem{theorem}{Theorem}
\title{Consequences of Matrices Sharing Eigenvalues and Eigenvectors}
\author{Kayla Miller, Steven J. Miller, Fuzhen Zhang}
\date{Source: arXiv:2609.10578v1 (CC BY 4.0)}
\begin{document}
\maketitle

\noindent\emph{Source and licence.} Consequences of Matrices Sharing Eigenvalues and Eigenvectors. Version arXiv:2609.10578v1 (CC BY 4.0), distributed by arXiv under the Creative Commons Attribution 4.0 International licence, http://creativecommons.org/licenses/by/4.0/, as recorded in the arXiv licence metadata for this exact version and shown on the abstract page https://arxiv.org/abs/2609.10578v1. Full licence notice: \texttt{licenses/arxiv-2609.10578-CC-BY-4.0.txt}.\par\smallskip

\noindent\textit{Editorial scope.} Counted unit: the article's theorem Thm:Normality (source0.tex lines 213--216). The article's complete original proof of it is delivered with it (source0.tex lines 588--590, one \texttt{proof} environment of 3 lines, which the article places away from the statement). No mathematics is written, repaired or re-derived: the statement and the proof are the article's own lines, in the article's own notation, and no retained line is substituted or re-flowed.  No further source-local context is needed: every cross-reference of every retained line resolves inside the two delivered spans.  Nothing was omitted from any retained span, so the package carries no omission record.  No retained line contains a citation, so the wrapper carries no bibliography and no bibliography entry.  No retained line contains a figure environment, an image-inclusion command or drawing code, so no image is delivered and the package declares no asset dependency.  Editorial formatting only, in addition: an AMS \texttt{amsart} preamble that restates the article's own declarations and macros used by the retained lines, namely the environments \texttt{theorem} on separate counters, together with the standard packages those lines need.  The retained environments print in this document as Theorem 1 in order of appearance, whereas the article prints the same items with the numbering of its own full document.  The complete native source of the article as distributed in the arXiv e-print for this exact version is bundled unchanged as \texttt{sources/209933/source0.tex}.
\begin{theorem}\label{Thm:Normality}
Every eigenvector  of $A$ is an eigenvector (regardless of the eigenvalues) of $A^*$ if and only if $A$ is \emph{normal},
that is,  $A^*A=AA^*$.
\end{theorem}

\begin{proof}  By Theorem \ref{Thm:Normality}, $A$ is normal. Thus, $A$ possesses $n$     orthonormal eigenvectors, say
$v_1, \dots, v_n$,  which  form a basis for $\Bbb C^n$. We have
 $(A-A^*)v_k=Av_k-A^*v_k =0$, $k=1, 2, \dots, n$. So, $A-A^*=0$ and $A$ is Hermitian. \end{proof}

\end{document}
